\section{Sujet n\textsuperscript{o}\,24}
\noindent\begin{minipage}{0.5\linewidth}\footnotesize Office du Baccalauréat du Cameroun\end{minipage}%
\hfill\begin{minipage}{0.45\linewidth}\footnotesize\raggedleft Examen : \textbf{Baccalauréat} · Série : \textbf{C/E}\\
Épreuve : \textbf{Mathématiques} · Durée : \textbf{4\,h} · Coef. : \textbf{7}\end{minipage}
\par\smallskip{\color{onbo}\rule{\linewidth}{1pt}}

\subsection*{Partie A — Évaluation des ressources \hfill (15 points)}

\noindent\textbf{Exercice 1.} \hfill\textit{(3 points)}\\
La cimenterie de Figuil enregistre, sur $6$ années consécutives, la consommation
d'électricité $x$ (en GWh) et la production de ciment $y$ (en milliers de tonnes) :
\begin{center}\small
\begin{tabular}{c|cccccc}
\toprule
$x_i$ & $20$ & $24$ & $28$ & $30$ & $34$ & $36$\\
$y_i$ & $52$ & $61$ & $68$ & $74$ & $83$ & $86$\\
\bottomrule
\end{tabular}
\end{center}
\begin{enumerate}[label=\arabic*)]
\item Calculer les moyennes $\bar x$ et $\bar y$, puis la covariance $\mathrm{Cov}(x,y)$. \hfill\textit{(1,25 pt)}
\item Déterminer une équation de la droite de régression de $y$ en $x$ par la méthode des moindres carrés. \hfill\textit{(1 pt)}
\item Estimer la production attendue pour une consommation de $40$~GWh. \hfill\textit{(0,75 pt)}
\end{enumerate}

\smallskip
\noindent\textbf{Exercice 2.} \hfill\textit{(3 points)}\\
Pour réguler son stock de clinker, l'usine modélise la masse (en milliers de tonnes)
présente au mois $n$ par la suite $(u_n)$ définie par $u_0=10$ et, pour tout $n\in\N$,
$u_{n+1}=0{,}75\,u_n+6$.
\begin{enumerate}[label=\arabic*)]
\item Calculer $u_1$ et $u_2$. \hfill\textit{(0,5 pt)}
\item On pose $v_n=u_n-24$. Montrer que $(v_n)$ est géométrique ; préciser sa raison et $v_0$. \hfill\textit{(1 pt)}
\item Exprimer $v_n$ puis $u_n$ en fonction de $n$, et déterminer $\displaystyle\lim_{n\to+\infty}u_n$. \hfill\textit{(1,5 pt)}
\end{enumerate}

\smallskip
\noindent\textbf{Exercice 3.} \hfill\textit{(4 points)}\\
Le coût marginal de fonctionnement d'un four est modélisé sur $\R$ par la fonction
$f$ définie par $f(x)=(2-x)\mathrm{e}^{x}$, de courbe $(\mathcal{C})$ dans un repère orthonormé
(unité : \SI{2}{\centi\metre}).
\begin{enumerate}[label=\arabic*)]
\item Déterminer les limites de $f$ en $-\infty$ et $+\infty$. \hfill\textit{(0,75 pt)}
\item Montrer que $f'(x)=(1-x)\mathrm{e}^{x}$ et dresser le tableau de variations de $f$. \hfill\textit{(1,25 pt)}
\item À l'aide d'une intégration par parties, montrer que $\displaystyle\int_0^{1}f(x)\,\dd x=2\mathrm{e}-3$. \hfill\textit{(1 pt)}
\item En déduire, en \unit{\centi\metre\squared}, l'aire du domaine limité par $(\mathcal{C})$, l'axe des abscisses et les droites $x=0$ et $x=1$. \hfill\textit{(1 pt)}
\end{enumerate}

\smallskip
\noindent\textbf{Exercice 4.} \hfill\textit{(5 points)}\\
Le plan complexe est muni d'un repère orthonormé direct $(O\,;\,\vec u,\vec v)$. On note
$A$, $B$ les points d'affixes $z_A=2$ et $z_B=3+i$.
\begin{enumerate}[label=\arabic*)]
\item Résoudre dans $\C$ l'équation $z^2-4z+8=0$ ; on note $z_1$ la solution de partie imaginaire positive. \hfill\textit{(1 pt)}
\item Écrire $z_1$ sous forme exponentielle, puis calculer $z_1^{\,4}$. \hfill\textit{(1 pt)}
\item Soit $S$ la similitude directe d'écriture complexe $z'=(1+i)z-2$. Déterminer son rapport, son angle et l'affixe $\omega$ de son centre. \hfill\textit{(1,5 pt)}
\item Déterminer l'affixe $z_C$ de l'image $C$ du point $B$ par $S$, puis la nature du triangle $\Omega B C$ où $\Omega$ est le centre de $S$. \hfill\textit{(1,5 pt)}
\end{enumerate}

\subsection*{Partie B — Évaluation des compétences \hfill (5 points)}

\noindent\textit{Le directeur technique de la cimenterie de Figuil te consulte, en tant qu'élève
de Terminale~C, pour trois décisions. (a) Il veut construire un silo cylindrique surmonté
d'un toit conique négligeable, dont la tôle utilisée pour la paroi latérale et le fond est
fixée à $24\pi$~\unit{\metre\squared} ; il cherche le rayon qui rend le volume du cylindre maximal.
(b) La production mensuelle de clinker, de \SI{30}{} milliers de tonnes ce mois-ci, croît de
\SI{4}{\percent} par mois ; il veut savoir quand elle dépassera \SI{45}{} milliers de tonnes.
(c) À la sortie d'usine, \SI{6}{\percent} des sacs ont un défaut d'emballage ; un contrôleur
en prélève $5$ au hasard.}

\begin{enumerate}[label=\textbf{Tâche \arabic*.},leftmargin=2.4em]
\item Déterminer le rayon et la hauteur du silo cylindrique de \emph{volume maximal}, ainsi que ce volume. \hfill\textit{(1,5 pt)}
\item Déterminer le mois à partir duquel la production mensuelle dépassera \SI{45}{} milliers de tonnes. \hfill\textit{(1,5 pt)}
\item Déterminer la probabilité qu'\emph{au moins un} des $5$ sacs prélevés présente un défaut. \hfill\textit{(1,5 pt)}
\end{enumerate}
\par\smallskip{\footnotesize\itshape Présentation générale : 0,5 point.}

% ════════════════════════════════════════════════════
\corrige{du sujet 24}

\noindent\textbf{Partie A — Exercice 1.}
1) $\bar x=\frac{20+24+28+30+34+36}{6}=\frac{172}{6}=\frac{86}{3}\approx28{,}67$ ;
$\bar y=\frac{52+61+68+74+83+86}{6}=\frac{424}{6}=\frac{212}{3}\approx70{,}67$.
$\sum x_iy_i=20\cdot52+24\cdot61+28\cdot68+30\cdot74+34\cdot83+36\cdot86=1040+1464+1904+2220+2822+3096=12546$.
$\mathrm{Cov}(x,y)=\frac1n\sum x_iy_i-\bar x\bar y=\frac{12546}{6}-\frac{86}{3}\cdot\frac{212}{3}=2091-\frac{18232}{9}\approx2091-2025{,}78=65{,}22$.
2) $V(x)=\frac1n\sum x_i^2-\bar x^2$ ; $\sum x_i^2=400+576+784+900+1156+1296=5112$,
$V(x)=\frac{5112}{6}-\big(\frac{86}{3}\big)^2=852-\frac{7396}{9}\approx852-821{,}78=30{,}22$.
Pente $a=\frac{\mathrm{Cov}(x,y)}{V(x)}\approx\frac{65{,}22}{30{,}22}\approx2{,}158$ ;
$b=\bar y-a\bar x\approx70{,}67-2{,}158\times28{,}67\approx70{,}67-61{,}87=8{,}80$. D'où $y\approx2{,}16\,x+8{,}80$.
3) Pour $x=40$ : $y\approx2{,}158\times40+8{,}80\approx95{,}1$, soit environ $\mathbf{95}$ milliers de tonnes.

\noindent\textbf{Exercice 2.}
1) $u_1=0{,}75\times10+6=13{,}5$ ; $u_2=0{,}75\times13{,}5+6=16{,}125$.
2) $v_{n+1}=u_{n+1}-24=0{,}75u_n+6-24=0{,}75u_n-18=0{,}75(u_n-24)=0{,}75\,v_n$ :
$(v_n)$ est géométrique de raison $q=0{,}75$ et $v_0=u_0-24=-14$.
3) $v_n=-14\times0{,}75^{\,n}$, donc $u_n=24-14\times0{,}75^{\,n}$. Comme $0<0{,}75<1$, $0{,}75^{\,n}\to0$
et $\displaystyle\lim_{n\to+\infty}u_n=\mathbf{24}$.

\noindent\textbf{Exercice 3.}
1) En $-\infty$ : $\mathrm{e}^x\to0$ et $(2-x)\to+\infty$ ; par croissances comparées $f(x)\to0^-$
(plus précisément $x\mathrm{e}^x\to0$, donc $f(x)\to0$). En $+\infty$ : $(2-x)\to-\infty$ et $\mathrm{e}^x\to+\infty$, donc $f(x)\to-\infty$.
2) $f'(x)=-\mathrm{e}^x+(2-x)\mathrm{e}^x=(1-x)\mathrm{e}^x$ : $f'(x)>0$ sur $]-\infty,1[$, $f'(x)<0$ sur $]1,+\infty[$.
$f$ croît sur $]-\infty,1]$, décroît sur $[1,+\infty[$, maximum $f(1)=\mathrm{e}$.
3) IPP avec $u=2-x$, $u'=-1$, $v'=\mathrm{e}^x$, $v=\mathrm{e}^x$ :
$\int_0^1(2-x)\mathrm{e}^x\dd x=\big[(2-x)\mathrm{e}^x\big]_0^1+\int_0^1\mathrm{e}^x\dd x=(1\cdot\mathrm{e}-2)+\big[\mathrm{e}^x\big]_0^1=(\mathrm{e}-2)+(\mathrm{e}-1)=2\mathrm{e}-3$.
4) Sur $[0,1]$, $f(x)=(2-x)\mathrm{e}^x>0$. L'aire en unités d'aire vaut $2\mathrm{e}-3$.
Une unité d'aire $=2\times2=\SI{4}{\centi\metre\squared}$, donc $\mathcal A=4(2\mathrm{e}-3)=8\mathrm{e}-12\approx\SI{9,75}{\centi\metre\squared}$.

\noindent\textbf{Exercice 4.}
1) $\Delta=16-32=-16=(4i)^2$ : $z=\frac{4\pm4i}{2}=2\pm2i$. Donc $z_1=2+2i$.
2) $|z_1|=\sqrt{4+4}=2\sqrt2$, $\arg z_1=\frac\pi4$ : $z_1=2\sqrt2\,\mathrm{e}^{i\pi/4}$.
$z_1^{\,4}=(2\sqrt2)^4\mathrm{e}^{i\pi}=64\times(-1)=-64$.
3) $z'=(1+i)z-2$. Rapport $=|1+i|=\sqrt2$ ; angle $=\arg(1+i)=\frac\pi4$.
Centre : $\omega=(1+i)\omega-2\Rightarrow\omega(1-(1+i))=-2\Rightarrow-i\,\omega=-2\Rightarrow\omega=\frac{2}{i}=-2i$.
4) $z_C=(1+i)(3+i)-2=(3+i+3i+i^2)-2=(3+4i-1)-2=4i$. Avec $\Omega$ d'affixe $-2i$ :
$\Omega B=|z_B-\omega|=|3+i+2i|=|3+3i|=3\sqrt2$ ; $\Omega C=|z_C-\omega|=|4i+2i|=|6i|=6$.
Or $S$ multiplie les distances par $\sqrt2$ : $\Omega C=\sqrt2\,\Omega B=\sqrt2\times3\sqrt2=6$ \checkmark.
L'angle $(\overrightarrow{\Omega B},\overrightarrow{\Omega C})=\frac\pi4$. Triangle $\Omega BC$ : rectangle en $B$
(car $\Omega C^2=36$, $\Omega B^2=18$, $BC^2=|z_C-z_B|^2=|{-3+3i}|^2=18$, et $18+18=36$),
donc \textbf{rectangle isocèle} de sommet de l'angle droit $B$.

\smallskip
\noindent\textbf{Partie B — Situation.}
\emph{Tâche 1.} Aire de tôle (fond + paroi latérale) : $\pi r^2+2\pi r h=24\pi$, soit $r^2+2rh=24$,
d'où $h=\frac{24-r^2}{2r}$ ($0<r<\sqrt{24}$). Volume $V(r)=\pi r^2h=\pi r^2\cdot\frac{24-r^2}{2r}=\frac\pi2(24r-r^3)$.
$V'(r)=\frac\pi2(24-3r^2)=0\Rightarrow r^2=8\Rightarrow r=2\sqrt2$. Alors $h=\frac{24-8}{2\cdot2\sqrt2}=\frac{16}{4\sqrt2}=2\sqrt2$.
Volume maximal : $V=\frac\pi2(24\cdot2\sqrt2-(2\sqrt2)^3)=\frac\pi2(48\sqrt2-16\sqrt2)=\frac\pi2\cdot32\sqrt2=16\sqrt2\,\pi\approx\SI{71,1}{\metre\cubed}$.

\emph{Tâche 2.} Production du mois $n$ (avec $n=0$ ce mois) : $p_n=30\times1{,}04^{\,n}$. On résout $p_n>45$,
soit $1{,}04^{\,n}>1{,}5$, d'où $n>\frac{\ln1{,}5}{\ln1{,}04}\approx\frac{0{,}4055}{0{,}0392}\approx10{,}3$ :
dès le \textbf{11\textsuperscript{e} mois} ($p_{11}=30\times1{,}04^{11}\approx\SI{46,2}{}$ milliers de tonnes).

\emph{Tâche 3.} $P(\text{au moins un défaut})=1-(0{,}94)^5\approx1-0{,}7339=\mathbf{0{,}266}$, soit environ \SI{26,6}{\percent}.
